%%%%%%%%%%%%%%%%%%%%%%%%%%%%%%%%%%%%%%%%%%%%%%%%%%%%%%%%%%%%%%%%%%%%%%
% Lớp 11 | Hàm số mũ và hàm số lôgarít | [Giảm] Các phương pháp giải được giảm tải
% Dạng 1D6?5-5: Toán thực tế áp dụng phương trình mũ, lôgarít
%%%%%%%%%%%%%%%%%%%%%%%%%%%%%%%%%%%%%%%%%%%%%%%%%%%%%%%%%%%%%%%%%%%%%%


% nguồn: 2026-11-15_11_1D6_Mu_Logarit (nhập 14/11/2026)
\begin{bt}%[1D6H5-5]%[ma:77cb3dfc]
% nguon: Archive/SP_KhoLatex/SP_Đề Theo Bài-Khối 11/1-B20.tex #4021
%[Dự án 2025 - Đề cấu trúc mới của Bộ theo bài]
Nếu $D_0$ là chênh lệch nhiệt độ ban đầu giữa một vật $M$ và các vật xung quanh, và nếu các vật xung quanh có nhiệt độ $T_S$, thì nhiệt độ của vật $M$ tại thời điểm $t$ được mô hình hóa bởi hàm số: $T(t)=T_S+D_0\cdot\mathrm{e}^{-kt}$ $(1)$ (trong đó $k$ là hằng số dương phụ thuộc vào vật $M$). Một con gà tây nướng được lấy từ lò nướng khi nhiệt độ của nó đã đạt đến $195^\circ$F và được đặt trên một bàn trong một căn phòng có nhiệt độ là $65^\circ$F. Nếu nhiệt độ của gà tây là $150^\circ$F sau nửa giờ, nhiệt độ của nó sau $60$ phút là bao nhiêu độ F? (làm tròn đến hàng đơn vị).
\shortans{$121$}
\loigiai
{
Ta có $T_S=65$ và độ chênh lệch nhiệt độ là $D_0=195-65=130$\\
Sau nửa giờ $(t=0{,}5)$ thì nhiệt độ của gà là $T=150$.\\
Áp dụng công thức $(1)\colon 150=65+130\cdot\mathrm{e}^{-k(0{,}5)} \Leftrightarrow \mathrm{e}^{-k}=\left(\dfrac{17}{26}\right)^2$.\\
Vậy $T(t)=65+130.\left(\dfrac{17}{26}\right)^{2t}$.\\
Suy ra nhiệt độ của gà sau $60$ phút ($t=1$ giờ) là $65+130\cdot\left(\dfrac{17}{26}\right)^{2\cdot1}\approx 121^\circ$F.
}
\end{bt}

% nguồn: 2026-11-15_11_1D6_Mu_Logarit (nhập 14/11/2026)
\begin{bt}%[1D6H5-5]%[ma:45e496f9]
% nguon: Archive/SP_KhoLatex/SP_Đề Theo Bài-Khối 11/1-B20.tex #4022
%[Dự án 2025 - Đề cấu trúc mới của Bộ theo bài]
Số lượng của một loài vi khuẩn trong phòng thí nghiệm được tính theo công thức $S(t)=A\mathrm{e}^{rt}$, trong đó $A$ là số lượng vi khuẩn ban đầu, $S(t)$ là số lượng vi khuẩn có sau $t$ (phút), $r$ là tỉ lệ tăng trưởng $(r>0)$, $t$ (tính theo phút) là thời gian tăng trưởng. Biết rằng số lượng vi khuẩn ban đầu có $500$ con và sau $6$ giờ có $2000$ con. Hỏi ít nhất bao nhiêu phút, kể từ lúc bắt đầu, số lượng vi khuẩn đạt ít nhất $120000$ con? (làm tròn đến hàng đơn vị).
\shortans{$1424$}
\loigiai
{
Ta có $A=500$, $6$ giờ = $360$ phút, $S(360)=2000$.\\
Sau $6$ giờ lượng vi khuẩn là $2000$ con, tức là
$$2000=500\mathrm{e}^{360r} \Leftrightarrow \mathrm{e}^{360r}=4 \Leftrightarrow r=\dfrac{\ln 4}{360} \ (\text{do} \ \mathrm{e}>1).$$
Số lượng vi khuẩn đạt ít nhất $120\,000$ con, nghĩa là
$$500\mathrm{e}^{\frac{\ln 4}{360}t}\geqslant 120\,000 \Leftrightarrow \dfrac{\ln 4}{360}\cdot t\geqslant 240 \Leftrightarrow \dfrac{\ln 4}{360}\cdot t \geqslant 240\Leftrightarrow t\geqslant \dfrac{360\ln 240}{\ln 4}\approx 1423{,}24 \ (\text{phút}).$$
Vậy sau ít nhất $1424$ phút thì số lượng vi khuẩn đạt ít nhất $120\,000$ con.
}
\end{bt}

% nguồn: 2026-11-15_11_1D6_Mu_Logarit (nhập 14/11/2026)
\begin{bt}%[1D6H5-5]%[ma:fa7fae34]
% nguon: Archive/SP_KhoLatex/SP_Đề Theo Bài-Khối 11/1-B20.tex #4026
%[Dự án 2025 - Đề cấu trúc mới của Bộ theo bài]
Dân số thế giới được tính theo công thức $S=A\mathrm{e}^{nr}$ trong đó $A$ là dân số của năm lấy làm mốc tính, $S$ là dân số sau $n$ năm, $r$ là tỉ lệ tăng dân số hằng năm. Cho biết năm $2005$ Việt Nam có khoảng $80\,902\,400$ người và tỉ lệ tăng dân số là $1{,}47\%$ một năm. Như vậy, nếu tỉ lệ tăng dân số hàng năm không đổi thì tối thiểu đến năm bao nhiêu dân của Việt Nam có khoảng $93\,713\,000$ người?
\shortans{$2015$}
\loigiai
{
Ta có $S=A\mathrm{e}^{nr}$, mốc năm $2005$ với số dân $80\,902\,400$ người, tỉ lệ $r=1{,}47\%=0{,}0147$ một năm.\\
Sau tối thiểu $n$ năm sau $2005$ thì số dân Việt Nam là $93\,713\,000$ tức là
$$93\,713\,000\leqslant 80\,902\,400\mathrm{e}^{0{,}01407n} \Leftrightarrow n\geqslant \dfrac{1}{0{,}0147}\cdot\ln \dfrac{93\,713\,000}{80\,902\,400}\approx 10.$$
Vậy sau ít nhất $10$ năm kể từ $2005$ hay tối thiểu đến năm $2005+10=2015$ thì dân số Việt Nam có khoảng $93\,713\,000$ người.
}
\end{bt}
